\documentclass{article}
\usepackage[margin=25mm]{geometry}
\usepackage{booktabs}
\begin{document}
\title{MC graph fragmentation diagnostic}\date{28 October 2026}\maketitle
The first 100 double-beta events and 100 electron events in the reduced
HDF5 catalogues were inspected. These are consecutive events, not a random
ensemble estimate. Source coordinates and all MC particles were used, without
diffusion or a fiducial cut. Voxel width is fixed at 0.5 mm. Each connected
component is counted, including single-node components. Radius equality is
accepted. Thresholds apply inclusively to summed voxel energy; kNN uses $k=10$
with the stated radius cap. Calculation uses Float32.

\begin{center}
\begin{tabular}{lrrrr}\toprule
Method & Radius [mm] & Threshold [keV] & $0\nu\beta\beta$: single & Electron: single\\\midrule
KDT & 1.25 & 1 & 1 & 1 \\
KDT & 1.25 & 0 & 1 & 1 \\
KDT & 2 & 1 & 45 & 44 \\
KDT & 3 & 1 & 57 & 46 \\
KDT & 5 & 1 & 60 & 49 \\
kNN & 1.25 & 1 & 1 & 1 \\
\bottomrule\end{tabular}
\end{center}
At the viewer defaults, the median component counts are 6.0 and 7.0,
respectively. Original Petit Float64 voxelisation and graph construction were
run independently on the same hits. Equality of complete component membership
for all tested events: double beta True; electrons True.

The strong radius dependence identifies short-distance graph fragmentation.
Removing the threshold does not recover the single-component population.
Switching to kNN with the same radius cap does not cure the observed splitting.
Petit uses these MC defaults and subsequently walks only its highest-energy
component. The viewer displays all components. Increasing the radius is a
diagnostic, not an adopted selection prescription. Remaining disconnected
deposits require event inspection and MC-truth validation. Viewer settings
remain unchanged.

Numerical source: \texttt{results/event\_analysis/mc\_graph\_diagnostic/results.h5}.
Original code: Petit \texttt{src/voxels.jl}, \texttt{src/track\_building.jl};
MC caller: \texttt{scripts/itaca\_single\_track\_analysis.jl}.
\end{document}
